\begin{figure}[!htbp]
  \centering
  \begin{tikzpicture}[
    x=1cm, y=1cm,
    >=Stealth,
    every node/.style={font=\small, align=center},
    arr/.style={-stealth, line width=0.95pt},
    qarr/.style={-stealth, line width=0.95pt, red!70!black, dashed},
    fp/.style={
      rectangle, rounded corners=3.5pt,
      draw=blue!58!black, line width=0.9pt,
      fill=blue!5, minimum width=2.75cm, minimum height=0.82cm
    },
    quant/.style={
      rectangle, rounded corners=3.5pt,
      draw=red!62!black, line width=0.95pt,
      fill=red!5, minimum width=2.35cm, minimum height=0.82cm
    },
    kpath/.style={
      rectangle, rounded corners=4pt,
      draw=cyan!55!black, line width=1.0pt,
      fill=cyan!8, minimum width=4.15cm, minimum height=1.06cm
    },
    vpath/.style={
      rectangle, rounded corners=4pt,
      draw=orange!78!black, line width=1.0pt,
      fill=orange!10, minimum width=4.15cm, minimum height=1.06cm
    },
    eqcard/.style={
      rectangle, rounded corners=5pt,
      draw=gray!55, line width=0.95pt,
      fill=gray!5, minimum width=12.95cm, minimum height=1.22cm
    },
    note/.style={
      rectangle, rounded corners=4pt,
      draw=gray!45, line width=0.9pt,
      fill=yellow!9, minimum width=9.9cm, minimum height=0.66cm
    }
  ]

    % Column titles
    \node[font=\bfseries\small, text=blue!68!black]   at (-4.95,3.06) {FP16 注意力流程};
    \node[font=\bfseries\small, text=red!72!black]    at (-0.42,3.06) {量化模块};
    \node[font=\bfseries\small, text=orange!82!black] at (4.55,3.06) {耦合传播路径};

    % Left column
    \node[fp] (in)     at (-4.95,2.25) {\((q, K, V)\)};
    \node[fp] (logits) at (-4.95,1.20) {\(z = qK^\top / \sqrt{d_k}\)};
    \node[fp] (attn)   at (-4.95,0.15) {\(a = \mathrm{softmax}(z)\)};
    \node[fp] (out)    at (-4.95,-0.90) {\(o = aV\)};

    \draw[arr, blue!68!black] (in.south) -- (logits.north);
    \draw[arr, blue!68!black] (logits.south) -- (attn.north);
    \draw[arr, blue!68!black] (attn.south) -- (out.north);

    % Middle column
    \node[font=\bfseries, text=red!72!black] at (-0.42,2.37) {\(Q(\cdot)\)};
    \node[quant] (kq) at (-0.42,1.45) {\(K \rightarrow \hat K\)};
    \node[quant] (vq) at (-0.42,-0.10) {\(V \rightarrow \hat V\)};

    \draw[qarr] ([xshift=2pt,yshift=5pt]in.east)  to[bend left=11]  (kq.west);
    \draw[qarr] ([xshift=2pt,yshift=-5pt]in.east) to[bend right=11] (vq.west);

    % Right column
    \node[kpath] (kp) at (4.55,1.55) {
      \textbf{Key 误差 \(\rightarrow\) logits 偏移}\\[-1pt]
      \(\rightarrow\) 分布扭曲 \(\rightarrow \hat a-a\)
    };
    \node[vpath] (vp) at (4.55,-0.20) {
      \textbf{Value 误差 \(\rightarrow\) 聚合内容偏移}\\[-1pt]
      \(\rightarrow\) 聚合侧传播项
    };

    \draw[arr, cyan!60!black]   (kq.east) -- (kp.west);
    \draw[arr, orange!82!black] (vq.east) -- (vp.west);

    % Bottom decomposition
    \node[eqcard] (eq) at (0.15,-2.12) {};
    \node[font=\footnotesize] at (0.15,-2.12) {%
      \textbf{输出误差分解：}\quad
      $\textstyle
      \hat o-o
      =
      \big[\sum_{i=1}^{S}\hat a_i(\hat v_i-v_i)\big]_{\text{Value 聚合项}}
      +
      \big[\sum_{i=1}^{S}(\hat a_i-a_i)v_i\big]_{\text{Key 分布项}}
      $%
    };

    % subtle guide marks
    \draw[blue!18, line width=0.6pt] (-3.15,2.75) -- (-3.15,-1.25);
    \draw[red!18, line width=0.6pt] (1.05,2.75) -- (1.05,-1.25);

  \end{tikzpicture}
  \caption{注意力误差的两条传播路径及其代数分解。FP16 attention 给出参考分布 \(a\) 与输出 \(o\)，量化后的 Key 误差经 logits 与 softmax 改变注意力分布，Value 误差在量化后权重下进入加权聚合。底部公式给出式~\eqref{eq:ch3-error-decomp} 的最终分解形式。}
  \label{fig:error-decomposition}
\end{figure}
